\section{Roadmap}
\label{sec:roadmap}

The package decomposition of \cref{sec:design-compart} fixes the shape;
what remains is the order. The sibling \code{ves-proof} repository
hosts the five Rust crates under \code{crates/};
\code{ves-proof-lean} hosts the five independent Lake packages under
\code{packages/}.  Their base skeletons now exist. \dotnetrs{} gains
scripts, two dev-dependencies, one CI job, and \code{TRUSTED.toml}.
Phases, each independently
valuable and abandonable without stranding the previous:

\begin{description}
\item[Phase 0 --- Hygiene (days).] \textbf{Complete} (commit
  \code{208a6c8b}, 2026-08-01). Fix the two live defects
  (\cref{sec:system-defects}): real alignment guards (or caller
  contracts) on \code{StackValue::load\_atomic}/\code{store\_atomic},
  and either a checked discipline or a narrower structure for the
  \code{SharedGlobalState} auto-trait overrides. Decide which dynamic
  gates can be promoted from \code{continue-on-error} to blocking (even
  a single Miri leg on \crate{dotnet-value} would be the first
  \emph{enforced} UB gate). None of this needs the DSL; all of it
  raises the floor the DSL builds on. All three line items landed in
  one commit: caller-contract alignment checks via \code{debug\_assert!}
  replace the vacuous hook; \code{SharedGlobalState} is now \code{!Send}
  by ordinary auto-trait inference instead of an asserted manual impl,
  with a minimal \code{AbortSignal} carrying the one fact that
  legitimately crosses the executor/waiter boundary; \code{miri-value}
  and \code{fuzz-raw-memory-access} are blocking CI legs (the broader
  Miri matrix and the other three fuzz targets remain
  \code{continue-on-error}, now in separate scheduled workflows); and
  the differential harness grew from one fixture
  (\code{expression\_compile\_42}) to seven, with a documented list of
  rejected candidates where \dotnetrs{} and real .NET diverge for
  reasons unrelated to this study (exit-code semantics for unhandled
  exceptions, an intentional alignment deviation, GC-timing-sensitive
  output). Phase 1 is now the critical path.
\item[Phase 1 --- The embedded layer (4--7 weeks).] \textbf{In
  progress.} In order:
  in \path{ves-proof/crates/}, \code{ves-syntax} (\textbf{done},
  commit \code{9082659}: script grammar, AST, canonical serializer,
  statement/subject hashing, golden corpus) and the declarative
  vocabulary v0 (\textbf{done}, commit \code{ad67707}: 77 F1--F9
  constructors across 16 namespaces with trust and family annotations,
  dependency/structural validation, and dual code generators emitting
  a \code{no\_std} Rust ghost-token module and a Lean declaration
  inventory, both pinned by golden-file tests);
  \code{ves-tokens} with cfg-mirrored constructors and
  \code{TrustGrant}s generated from \code{TRUSTED.toml} (seeded from
  the accepted-risk boundaries); \code{ves-macros} with
  token-identity snapshot tests; \code{ves-check} in structural mode
  (empty proof manifest, everything trusted or unproved); scripts on
  every unsafe block in \crate{dotnet-value}, \crate{dotnet-utils},
  and \crate{dotnet-runtime-memory}, and \code{ves::contract!}
  declarations on the 62 documented \code{unsafe fn}. Deliverable:
  premise existence, cfg-liveness, and premise dataflow checked by
  rustc in every CI matrix leg --- defect (a) becomes uncompilable
  (\cref{sec:dsl-example-fail}) --- plus structural coverage
  enforcement. This is a state-of-the-art contribution at (beyond) the
  tag-std level~\cite{tagstd2025} before any theorem exists, and it is
  also the test of the central hypothesis: if scripts cannot survive
  this project's refactor velocity, the lesson costs a linter rather
  than a model.
\item[Phase 2 --- The model (3--6 months).] \code{VesModel}:
  in \path{ves-proof-lean/packages/VesModel/},
  $\mathcal{M}_{\mathrm{VES}}$ in Lean 4 --- layout rules, pointer
  taxonomy, memory/atomicity model, eval-stack typing,
  reachability/finalization; executable; differentially tested against
  real .NET on the fixture corpus (which also expands the diff harness
  the project needs anyway). Independent of Phase 1's crates by
  design. Deliverable: the first mechanized ECMA-335 core ---
  publishable, and useful as a specification oracle even if no later
  phase completes.
\item[Phase 3 --- The two hard theories (2--4 months, parallelizable).]
  Under \path{ves-proof-lean/packages/}, \code{VesProtocol}: the
  STW/arena LTS, model-checked exhaustively at
  small configurations, then its invariants proved.
  \code{RustAssumptions}: the $\mathcal{M}_{\mathrm{Rust}}$
  axiomatization with Miri-falsifier templates. \code{VesCore}'s goal
  elaborator and stub generator land here, making the manifests flow
  end-to-end. Gate: the protocol theorem
  (\cref{sec:risks-concurrency}) --- if it stalls, F1 sites stay on
  trust grants and the campaign continues at reduced strength.
\item[Phase 4 --- The campaign (3--6 months, then open-ended).]
  In \path{ves-proof-lean/packages/DotnetRsProofs/}: the
  representation relation $\mathcal{R}$;
  \code{GcDescFaithful}; then the critical sites in value order --- the two
  accepted-risk boundaries, the \code{Collect} impls, the atomics
  family, the P/Invoke pinning contracts. Each proved lemma retires a
  trust grant or an unproved hash under the ratchet, a one-line script
  diff per site (\cref{sec:dsl-example-rebrand}). Stretch goals beyond
  the stated scope, in order of leverage: MIR-level tethering for the
  critical modules; a mechanized CIL verifier fragment to discharge the
  managed-IL trust class; RefinedRust/Rocq treatment of the \gcarena{}
  boundary.
\end{description}

Two standing rules constrain the effort: every phase must deliver an
artifact \dotnetrs{} uses in CI within that phase, and the
trusted-premise ratchet only moves downward.
